La irrupción de la inteligencia artificial generativa en el ámbito educativo plantea una paradoja, las herramientas son capaces de resolver en segundos los ejercicios que los estudiantes de programación deben aprender a resolver por sí mismos. Prohibirlas resulta contraproducente en un contexto en el que su uso profesional es ya inevitable, y permitirlas sin control conduce al aprendizaje pasivo y a la dependencia cognitiva.

Este Trabajo Fin de Grado presenta \textbf{SocratiCode}, una plataforma web de tutoría de programación que propone una tercera vía. En lugar de generar la solución, el sistema emplea un modelo de lenguaje local instruido mediante \textit{prompt engineering} para actuar como tutor socrático. Este, guía al alumno con preguntas que le llevan a descubrir el error y la solución por sí mismo, y se niega explícitamente a proporcionar código corregido bajo cualquier circunstancia.

La plataforma integra un editor de código con ejecución en un entorno aislado, un canal de comunicación en tiempo real basado en \textit{Server-Sent Events} y un sistema de moderación dual que protege el entorno educativo. Al operar íntegramente en local, no depende de APIs comerciales ni expone datos académicos a terceros.

La interfaz, diseñada como una aplicación de página única, integra en una misma pantalla el editor de código, la terminal de salida y el chat con el tutor. Esta disposición permite al alumno ver en un solo vistazo el error que ha cometido, la salida que ha producido y la pista socrática que el tutor le ofrece, eliminando la necesidad de alternar entre aplicaciones durante el proceso de aprendizaje.

\palabrasclave{Modelos de lenguaje grandes (LLM), ingeniería de prompts, inteligencia artificial generativa, tutor socrático, enseñanza de programación, ejecución de código segura}
